\documentclass{amsart}
% For dealing with references we use the comment environment
\usepackage{verbatim}
\newenvironment{reference}{\comment}{\endcomment}
%\newenvironment{reference}{}{}
\newenvironment{slogan}{\comment}{\endcomment}
\newenvironment{history}{\comment}{\endcomment}

% For commutative diagrams we use Xy-pic
\usepackage[all]{xy}

% We use 2cell for 2-commutative diagrams.
\xyoption{2cell}
\UseAllTwocells

% We use multicol for the list of chapters between chapters
\usepackage{multicol}

% This is generally recommended for better output
\usepackage{lmodern}
\usepackage[T1]{fontenc}

% For cross-file-references

% Package for hypertext links:
\usepackage{hyperref}

% For any local file, say "hello.tex" you want to link to please

% Theorem environments.
%
\theoremstyle{plain}
\newtheorem{theorem}[subsection]{Theorem}
\newtheorem{proposition}[subsection]{Proposition}
\newtheorem{lemma}[subsection]{Lemma}

\theoremstyle{definition}
\newtheorem{definition}[subsection]{Definition}
\newtheorem{example}[subsection]{Example}
\newtheorem{exercise}[subsection]{Exercise}
\newtheorem{situation}[subsection]{Situation}

\theoremstyle{remark}
\newtheorem{remark}[subsection]{Remark}
\newtheorem{remarks}[subsection]{Remarks}

\numberwithin{equation}{subsection}

% Macros
%
\def\lim{\mathop{\mathrm{lim}}\nolimits}
\def\colim{\mathop{\mathrm{colim}}\nolimits}
\def\Spec{\mathop{\mathrm{Spec}}}
\def\Hom{\mathop{\mathrm{Hom}}\nolimits}
\def\Ext{\mathop{\mathrm{Ext}}\nolimits}
\def\SheafHom{\mathop{\mathcal{H}\!\mathit{om}}\nolimits}
\def\SheafExt{\mathop{\mathcal{E}\!\mathit{xt}}\nolimits}
\def\Sch{\mathit{Sch}}
\def\Mor{\mathop{\mathrm{Mor}}\nolimits}
\def\Ob{\mathop{\mathrm{Ob}}\nolimits}
\def\Sh{\mathop{\mathit{Sh}}\nolimits}
\def\NL{\mathop{N\!L}\nolimits}
\def\CH{\mathop{\mathrm{CH}}\nolimits}
\def\proetale{{pro\text{-}\acute{e}tale}}
\def\etale{{\acute{e}tale}}
\def\QCoh{\mathit{QCoh}}
\def\Ker{\mathop{\mathrm{Ker}}}
\def\Im{\mathop{\mathrm{Im}}}
\def\Coker{\mathop{\mathrm{Coker}}}
\def\Coim{\mathop{\mathrm{Coim}}}

% Boxtimes
%
\DeclareMathSymbol{\boxtimes}{\mathbin}{AMSa}{"02}

%
% Macros for moduli stacks/spaces
%
\def\QCohstack{\mathcal{QC}\!\mathit{oh}}
\def\Cohstack{\mathcal{C}\!\mathit{oh}}
\def\Spacesstack{\mathcal{S}\!\mathit{paces}}
\def\Quotfunctor{\mathrm{Quot}}
\def\Hilbfunctor{\mathrm{Hilb}}
\def\Curvesstack{\mathcal{C}\!\mathit{urves}}
\def\Polarizedstack{\mathcal{P}\!\mathit{olarized}}
\def\Complexesstack{\mathcal{C}\!\mathit{omplexes}}
% \Pic is the operator that assigns to X its picard group, usage \Pic(X)
% \Picardstack_{X/B} denotes the Picard stack of X over B
% \Picardfunctor_{X/B} denotes the Picard functor of X over B
\def\Pic{\mathop{\mathrm{Pic}}\nolimits}
\def\Picardstack{\mathcal{P}\!\mathit{ic}}
\def\Picardfunctor{\mathrm{Pic}}
\def\Deformationcategory{\mathcal{D}\!\mathit{ef}}

\setlength{\emergencystretch}{3em}
\begin{document}
\title[Stacks Project, Tag 0BRK]{419 --- Galois action on primes and residue automorphisms in infinite extensions}
\maketitle
\noindent Copyright (C) 2005--2025 Johan de Jong. Stacks Project authors. Source: \href{https://stacks.math.columbia.edu/tag/0BRK}{Tag 0BRK}.

\noindent Editorial difficulty note (not author text). Difficulty H2: The proof uses a maximal compatible pair of a Galois subfield and an automorphism moving a prescribed prime. Finite normal extensions and correction by a stabilizer force enlargement; a second construction simultaneously matches the prescribed residue-field automorphism..

\noindent Editorial provenance note (not author text). History: excerpt from revision \texttt{a0444\allowbreak 6e57e\allowbreak c1fbc\allowbreak 252a8\allowbreak 71afc\allowbreak ec775\allowbreak 2fb28\allowbreak 07b14}, more-algebra.tex, lines 32776--32866. Reference presentation was adapted; original-author context and core support blocks were retained or added. The mathematical wording is unchanged. Licensed under GNU FDL 1.2 or later; no invariant sections or cover texts. See ../licenses/COPYING. Context blocks reproduce author definitions and supporting results; they are not additional counted problems.

\begin{lemma}
\label{lemma-one-orbit-geometric-galois}
Let $A$ be a normal domain with fraction field $K$.
Let $L/K$ be a (possibly infinite) Galois extension.
Let $G = \text{Gal}(L/K)$ and let
$B$ be the integral closure of $A$ in $L$.
\begin{enumerate}
\item For any two primes
$\mathfrak q, \mathfrak q' \subset B$ lying over the same prime in $A$
there exists a $\sigma \in G$ with $\sigma(\mathfrak q) = \mathfrak q'$.
\item Let $\mathfrak q \subset B$ be a prime lying over
$\mathfrak p \subset A$. Then $\kappa(\mathfrak q)/\kappa(\mathfrak p)$
is an algebraic normal extension and the map
$$
D = \{\sigma \in G \mid \sigma(\mathfrak q) = \mathfrak q\}
\longrightarrow
\text{Aut}(\kappa(\mathfrak q)/\kappa(\mathfrak p))
$$
is surjective.
\end{enumerate}
\end{lemma}

\begin{proof}
Proof of (1). Consider pairs $(M, \sigma)$ where $K \subset M \subset L$
is a subfield such that $M/K$ is Galois, $\sigma \in \text{Gal}(M/K)$
with $\sigma(\mathfrak q \cap M) = \mathfrak q' \cap M$.
We say $(M', \sigma') \geq (M, \sigma)$ if and only if
$M \subset M'$ and $\sigma'|_M = \sigma$.
Observe that $(K, \text{id}_K)$ is such a pair as $A = K \cap B$
since $A$ is a normal domain.
The collection of these pairs satisfies the hypotheses of Zorn's lemma,
hence there exists a maximal pair $(M, \sigma)$.
If $M \not = L$, then we can find
$M \subset M' \subset L$ with $M'/M$ nontrivial and finite and $M'/K$ Galois
(Fields, Lemma \href{https://stacks.math.columbia.edu/tag/0BMG}{lemma normal closure inside normal (Tag 0BMG)}).
Choose $\sigma' \in \text{Gal}(M'/K)$ whose restriction to $M$
is $\sigma$ (Fields, Lemma \href{https://stacks.math.columbia.edu/tag/0BMK}{lemma galois infinite (Tag 0BMK)}).
Then the primes $\sigma'(\mathfrak q \cap M')$ and $\mathfrak q' \cap M'$
restrict to the same prime of $B \cap M$. Since
$B \cap M = (B \cap M')^{\text{Gal}(M'/M)}$ we can
use Lemma \href{https://stacks.math.columbia.edu/tag/0BRI}{lemma one orbit (Tag 0BRI)} to find $\tau \in \text{Gal}(M'/M)$
with $\tau(\sigma'(\mathfrak q \cap M')) = \mathfrak q' \cap M'$.
Hence $(M', \tau \circ \sigma') > (M, \sigma)$
contradicting the maximality of $(M, \sigma)$.

\medskip\noindent
Part (2) is proved in exactly the same manner as part (1). We
write out the details. Pick
$\overline{\sigma} \in \text{Aut}(\kappa(\mathfrak q)/\kappa(\mathfrak p))$.
Consider pairs $(M, \sigma)$ where $K \subset M \subset L$
is a subfield such that $M/K$ is Galois, $\sigma \in \text{Gal}(M/K)$
with $\sigma(\mathfrak q \cap M) = \mathfrak q \cap M$ and
$$
\xymatrix{
\kappa(\mathfrak q \cap M) \ar[r] \ar[d]_\sigma &
\kappa(\mathfrak q) \ar[d]_{\overline{\sigma}} \\
\kappa(\mathfrak q \cap M) \ar[r] & \kappa(\mathfrak q)
}
$$
commutes. We say $(M', \sigma') \geq (M, \sigma)$ if and only if
$M \subset M'$ and $\sigma'|_M = \sigma$.
As above $(K, \text{id}_K)$ is such a pair.
The collection of these pairs satisfies the hypotheses of Zorn's lemma,
hence there exists a maximal pair $(M, \sigma)$.
If $M \not = L$, then we can find
$M \subset M' \subset L$ with $M'/M$ finite and $M'/K$ Galois
(Fields, Lemma \href{https://stacks.math.columbia.edu/tag/0BMG}{lemma normal closure inside normal (Tag 0BMG)}).
Choose $\sigma' \in \text{Gal}(M'/K)$ whose restriction to $M$
is $\sigma$ (Fields, Lemma \href{https://stacks.math.columbia.edu/tag/0BMK}{lemma galois infinite (Tag 0BMK)}).
Then the primes $\sigma'(\mathfrak q \cap M')$ and $\mathfrak q \cap M'$
restrict to the same prime of $B \cap M$. Adjusting the choice
of $\sigma'$ as in the first paragraph, we may assume that
$\sigma'(\mathfrak q \cap M') = \mathfrak q \cap M'$.
Then $\sigma'$ and $\overline{\sigma}$ define maps
$\kappa(\mathfrak q \cap M') \to \kappa(\mathfrak q)$
which agree on $\kappa(\mathfrak q \cap M)$. Since
$B \cap M = (B \cap M')^{\text{Gal}(M'/M)}$ we can
use Lemma \href{https://stacks.math.columbia.edu/tag/0BRJ}{lemma one orbit geometric (Tag 0BRJ)}
to find $\tau \in \text{Gal}(M'/M)$ with
$\tau(\mathfrak q \cap M') = \mathfrak q \cap M'$
such that $\tau \circ \sigma$ and $\overline{\sigma}$
induce the same map on $\kappa(\mathfrak q \cap M')$.
There is a small detail here in that the lemma first
guarantees that $\kappa(\mathfrak q \cap M')/\kappa(\mathfrak q \cap M)$
is normal, which then tells us that the difference between
the maps is an automorphism of this extension
(Fields, Lemma \href{https://stacks.math.columbia.edu/tag/0BR4}{lemma normal embeddings differ by aut (Tag 0BR4)}),
to which we can
apply the lemma to get $\tau$. Hence $(M', \tau \circ \sigma') > (M, \sigma)$
contradicting the maximality of $(M, \sigma)$.
\end{proof}
\end{document}
